\section{A Gamma integral basis and eventual reconstruction at every fixed gap}
\label{gb:section}

The preceding sections develop identities by localizing a finite-difference
operator at its boundary.  We now apply a different analytic construction
to a remaining reconstruction problem in the repository: the optimal
initial moments of two unequal even harmonic tails.  The inherited
fixed-gap reduction expresses the obstruction as a rational determinant
of size half the gap \cite{repo-gamma}.  We construct its entire null
space by ordinary Gamma integrals, determine the exact normalization,
and obtain a nonzero determinant asymptotic.  Consequently the optimal
prefix holds eventually on \emph{every fixed-gap ray}.  This differs from
the previously established eventual result with the smaller order fixed
\cite{repo:quartic}:
here both orders tend to infinity while their difference remains fixed.

The Gamma recurrence, Stirling expansion, and Fourier-transform setting
are classical \cite{dlmf:gamma,dlmf-stirling,ramanujan-transform}.
The contributions established below are the explicit basis for this
particular null equation, its exact bridge to the inherited rational
determinant, the fixed-gap asymptotic, and its harmonic-number correction.
This is a statement of the advance relative to the inspected material;
it does not assert historical priority.  Throughout the section, rank
means rank of the specified rational \emph{formal coordinate matrix}.
No arithmetic independence of evaluated zeta values or moments is
assumed or concluded.

\subsection{The inherited moments and the determinant to be studied}

For integers $1\le a<b$, write
\[
 p=2a,\qquad q=2b,\qquad m=b-a\ge1,
 \qquad x_n=n+\tfrac12.
\]
Thus $b$ is half the larger order and $m$ is half the gap.  Define
\begin{equation}\label{gb:tails}
 A_p(n)=\zeta(p,n+1)=\zeta(p)-H_n^{(p)},\qquad
 D_{p,q}(s)=\sum_{n=0}^{\infty}A_p(n)A_q(n)x_n^{-s}.
\end{equation}
The defining series converges absolutely for $\Rea s>3-p-q$.
Let $D_r=D_{p,q}(-2r)$, using meromorphic continuation when necessary.
These points are regular.  Indeed the centered Hurwitz expansion has
the form
\[
 A_p(n)\sim\frac{x_n^{1-p}}{p-1}
 +\sum_{k\ge1}\frac{B_{2k}(1/2)}{(2k)!}
       (p)_{2k-1}x_n^{1-p-2k}.
\]
Subtracting finitely many terms from the product continues
\eqref{gb:tails} successively to larger half-planes.  Its possible poles
lie at $s=3-p-q-2\ell$, $\ell\ge0$, all odd integers.  In particular
the even arguments used in $D_r$ are regular.

The inherited two-tail calculus supplies blocks $C_j$ and the mutually
inverse Bernoulli transforms
\begin{align}
 D_r&=\sum_{j=0}^r\frac{\binom{2r}{2j}}{2j+1}
          B_{2r-2j}(1/2)C_j,\label{gb:forward}\\
 C_j&=\sum_{r=0}^j\binom{2j+1}{2r}2^{-2(j-r)}D_r.
 \label{gb:inverse}
\end{align}
Here $B_n(x)$ denotes a Bernoulli polynomial and $B_n=B_n(0)$, with
$B_1=-1/2$.  For $j\ge b$ the late coordinates are
\begin{equation}\label{gb:late}
 C_j=\sum_{k=0}^{b-1}a_{jk}Z_k,\qquad
 Z_k=\zeta(2k),\qquad Z_0=-\tfrac12.
\end{equation}
Their exact rational coefficients, also inherited, are as follows.
Put $d=2j+1$ and $e=d-p-q+2k+1$.  If $e<0$ set $a_{jk}=0$;
otherwise set
\begin{equation}\label{gb:source-entry}
 a_{jk}=B_e\left\{
 \frac{\binom{d-q+1}{p-2k}}{d-q+1}
 +\frac{\binom{d-p+1}{q-2k}}{d-p+1}\right\}.
\end{equation}
Each binomial summand is omitted unless its lower index lies between
$2$ and its upper index.  In particular, a lower index of zero is
\emph{omitted}.  This convention is essential: it determines the
constant subtraction in the differential equation below.  Equations
\eqref{gb:forward}--\eqref{gb:source-entry} are inputs from
\cite[Section~2, inherited moment coordinates]{repo-gamma}.

Define the $b\times b$ late-block matrix
\begin{equation}\label{gb:M}
 M_{a,b}=(a_{b+r,k})_{0\le r,k<b}.
\end{equation}
On the ordered formal labels
\[
 C_0,\ldots,C_{b-1},\quad Z_0,\ldots,Z_{b-1},
\]
the coefficient matrix of $C_0,\ldots,C_{2b-1}$ is
$\operatorname{diag}(I_b,M_{a,b})$.  The forward transform
\eqref{gb:forward} is triangular with diagonal $1/(2r+1)$.
Consequently the determinant of the complete matrix for the first
$q=2b$ moments $D_0,\ldots,D_{2b-1}$ is
\begin{equation}\label{gb:full-prefix}
             \frac{\det M_{a,b}}{(4b-1)!!}.
\end{equation}
Thus the optimal-prefix question is exactly the nonvanishing of
\eqref{gb:M}.  If it is nonzero, those first $2b$ moments generate
every later moment in the inherited formal coordinate system, and
evaluation gives the corresponding rational generation statements for
the actual constants.  Fewer than $2b$ rows cannot have rank $2b$ on
these $2b$ specified labels.

\subsection{The fixed-gap recurrence and its exact normalization}

For the remainder of the section fix $m\ge1$.  To distinguish the
recurrence coefficients from both the tails and Bernoulli numbers, put
\begin{equation}\label{gb:recurrence-factors}
 \mathfrak a_r=\begin{cases}
 0,&1\le r\le m,\\
 (2r)^{\underline{2m}},&r>m,
 \end{cases}
 \qquad
 \mathfrak b_N=(2N-1)^{\underline{2m}},\qquad
 d_N=2(\mathfrak a_N+\mathfrak b_N)\quad(N>m).
\end{equation}
The underline denotes a falling factorial.  All $d_N$ in this range
are strictly positive.  Let
\begin{equation}\label{gb:beta}
 \beta_k=\frac{2|B_{2k}|}{(2k)!}
        =\frac{4\zeta(2k)}{(2\pi)^{2k}}>0,
 \qquad k\ge1.
\end{equation}
The rational fundamental array is defined by
\begin{align}
 C^{[m]}_{r,j}&=\delta_{rj}&& (1\le r,j\le m),\label{gb:initial}\\
 C^{[m]}_{N,j}&=\frac1{d_N}\sum_{r=1}^{N-1}
   \beta_{N-r}(\mathfrak a_r+\mathfrak b_N)C^{[m]}_{r,j}
       &&(N>m,\ 1\le j\le m).\label{gb:recurrence}
\end{align}
The superscript separates this array from the moment blocks $C_j$.
For $b>m$ write
\begin{equation}\label{gb:F}
                  F_{b,m}=(C^{[m]}_{b+i,j})_{1\le i,j\le m}.
\end{equation}

For an even formal series $R(z)=\sum_{n\ge1}r_nz^{2n}$, define
\begin{equation}\label{gb:cot-operator}
 \mathcal T_mR
 =\bigl(R^{(2m)}-R^{(2m)}(0)\bigr)\cot(z/2)
                      +\bigl(R\cot(z/2)\bigr)^{(2m)}.
\end{equation}
Both terms combine to an odd formal Taylor series: the subtraction
removes the possible pole in the first product, while $R\cot(z/2)$
is an odd Taylor series.  Expanding
\[
 \cot(z/2)=\frac2z-\sum_{k\ge1}\beta_kz^{2k-1}
\]
shows that, for $N>m$, its coefficient at $z^{2N-2m-1}$ is
\begin{equation}\label{gb:coefficient}
 d_Nr_N-\sum_{r=1}^{N-1}
      \beta_{N-r}(\mathfrak a_r+\mathfrak b_N)r_r.
\end{equation}
The value $\mathfrak a_m=0$ records precisely the subtracted
$R^{(2m)}(0)$.  Since $d_N>0$, the first $m$ coefficients are
arbitrary and determine all later ones uniquely.  The formal null
space of $\mathcal T_m$ has dimension exactly $m$.

\begin{proposition}[Inherited boundary reduction, with normalization]
\label{gb:bridge}
For $b>m\ge1$,
\begin{equation}\label{gb:bridge-formula}
 \boxed{\displaystyle
 \det M_{b-m,b}
 =\frac{(-1)^{b(b-1)/2}}{4^bb!}
       \left(\prod_{N=m+1}^{b+m}d_N\right)\det F_{b,m}.}
\end{equation}
This is the complementary-minor theorem of
\cite[Theorem~\texttt{gap:complementary}]{repo-gamma}.
\end{proposition}

\begin{proof}
We reproduce the normalization argument so that the subsequent
Gamma determinant can be checked against the original coordinates.
Set $h(t)=\coth(t/2)$ and
\[
 \mathcal L_mQ=(Q^{(2m)}-Q^{(2m)}(0))h+(Qh)^{(2m)}.
\]
Let $K_{m,b}$ map the coefficients $(q_1,\ldots,q_b)$ of
$Q(t)=\sum_{n=1}^bq_nt^{2n}$ to the coefficients of
$t,t^3,\ldots,t^{2b-1}$ in $\mathcal L_mQ$.
Write $h(t)=\sum_{k\ge0}h_kt^{2k-1}$, where $h_0=2$ and
$h_k=2B_{2k}/(2k)!$ for $k\ge1$.  In row $r$, indexed from zero,
put $N=r+m+1$.  The entry in column $n$ is
\[
 h_{N-n}(\mathfrak a_n+\mathfrak b_N)\quad(n\le N),
\]
and is zero for $n>N$.  In \eqref{gb:source-entry}, substitution
$j=b+r$, $k=b-n$ gives Bernoulli index $2(N-n)$ and the upper/lower
binomial pairs
\[
 (2r+2,\,2n-2m),\qquad(2r+2m+2,\,2n).
\]
Multiplication by $2(2n)!/(2r+1)!$ therefore gives exactly the
displayed operator entry, including the omission when $n\le m$.
If $M^{\mathrm{rev}}$ denotes column reversal, then
\[
 K_{m,b}=
 \operatorname{diag}_{0\le r<b}\!\left(\frac2{(2r+1)!}\right)
 M_{b-m,b}^{\mathrm{rev}}
 \operatorname{diag}_{1\le n\le b}\!\bigl((2n)!\bigr),
\]
so
\begin{equation}\label{gb:operator-conjugacy}
 \det K_{m,b}=(-1)^{b(b-1)/2}4^bb!\det M_{b-m,b}.
\end{equation}

The relation $h_k=(-1)^{k-1}\beta_k$ for $k\ge1$ and
\eqref{gb:recurrence} imply that
\[
 Q_j(t)=\sum_{n\ge1}(-1)^{n-j}C^{[m]}_{n,j}t^{2n},
 \qquad1\le j\le m,
\]
are formal null solutions of $\mathcal L_m$, with initial
$m\times m$ coefficient block equal to the identity.  Replace
the first $m$ columns of $K_{m,b}$ by the corresponding linear
combinations given by the truncations of these solutions through
$n=b$.  This changes no determinant.  The first $b-m$ rows of
those columns vanish.  Their bottom $m$ rows equal $-TQ_+$, where
$T$ is lower triangular with diagonal $d_{b+1},\ldots,d_{b+m}$ and
\[
 (Q_+)_{ij}=(-1)^{b+i-j}C^{[m]}_{b+i,j}.
\]
Indeed each truncated null equation is the negative of its omitted
future terms, and its newest future coefficient has factor $d_{b+i}$.
Moving the first $m$ columns to the end leaves the other diagonal
block triangular with diagonal $d_{m+1},\ldots,d_b$.  It follows that
\[
 \det K_{m,b}=(-1)^{m(b-m)+m}
       \left(\prod_{N=m+1}^{b+m}d_N\right)\det Q_+.
\]
Factoring its row and column signs gives
$\det Q_+=(-1)^{mb}\det F_{b,m}$.  The combined exponent
$2mb-m(m-1)$ is even.  Equation \eqref{gb:operator-conjugacy}
now proves \eqref{gb:bridge-formula}.
\end{proof}

\subsection{An explicit Gamma integral basis}

Define the positive even weight
\begin{equation}\label{gb:W}
 W_m(x)=\left|\Gamma\!\left(\tfrac14+\tfrac{ix}{2}\right)\right|^{4m}
 \qquad(x\in\mathbb R)
\end{equation}
and the odd multipliers
\begin{equation}\label{gb:P}
 P_j(x)=\sinh\!\bigl(\pi(m-j)x\bigr)\quad(1\le j<m),
 \qquad P_m(x)=\tanh(\pi x).
\end{equation}
Their order is chosen to agree with their distinct decay scales.

\begin{theorem}[Gamma basis and exact top moments]\label{gb:basis}
For $1\le j\le m$, the integral
\begin{equation}\label{gb:R}
 R_j(z)=4\sin(z/2)\int_0^\infty
                  \sinh(xz)P_j(x)W_m(x)\dd x
\end{equation}
defines an even holomorphic function in $|\Rea z|<j\pi$.
Its germ at zero satisfies $\mathcal T_mR_j=0$ and
\begin{equation}\label{gb:top-derivative}
 R_j^{(2m)}(0)=0\quad(j<m),\qquad
 R_m^{(2m)}(0)=(-1)^{m+1}2^{2m+2}\pi^m.
\end{equation}
The $m$ germs form a basis of the formal null space of
\eqref{gb:cot-operator}.  Their Taylor radii are exactly $j\pi$.
\end{theorem}

The contour argument below proves the null equation and the exact
derivatives.  The basis assertion will follow from the positive
initial determinant in \cref{gb:normalization}; the exact Taylor
radii will follow independently from \cref{gb:coefficient-asymptotic}.
This order of proof avoids assuming either linear independence or
the location of a singularity from the integral representation alone.

\begin{proof}[Contour proof]
Stirling's formula gives, as $x\to+\infty$,
\[
 W_m(x)=2^m(2\pi)^{2m}x^{-m}e^{-m\pi x}(1+O(x^{-2})).
\]
Together with \eqref{gb:P}, this proves absolute convergence in
the stated strips, locally uniform convergence of all differentiated
integrals, and evenness of \eqref{gb:R}.  We use the usual meromorphic
Gamma function and the integer-power continuation
\[
 W_m(y)=\left[
 \Gamma(\tfrac14+iy/2)\Gamma(\tfrac14-iy/2)\right]^{2m}.
\]
There is no branch choice for this continuation; the absolute-value
notation in \eqref{gb:W} is used only on the real line.

For any one of the allowed odd multipliers $P$, put
\[
 G_P(t)=\int_{\mathbb R}e^{ixt}P(x)W_m(x)\dd x,
 \qquad Q_P(t)=2\sinh(t/2)G_P(t).
\]
These integrals and all their derivatives converge absolutely for
real $t$ and define analytic germs.  Since $G_P$ is odd, $Q_P$ is
even and vanishes at zero.  With $D=d/dt$, let
\[
 E_P(t)=e^t(D+\tfrac12)^{2m}G_P(t)
                 -e^{-t}(D-\tfrac12)^{2m}G_P(t).
\]
The identities $Q_P=(e^{t/2}-e^{-t/2})G_P$ and
$Q_P\coth(t/2)=(e^{t/2}+e^{-t/2})G_P$ yield, by differentiation,
\begin{equation}\label{gb:contour-algebra}
 \coth(t/2)Q_P^{(2m)}(t)
       +(Q_P(t)\coth(t/2))^{(2m)}
                    =\frac{E_P(t)}{\sinh(t/2)}.
\end{equation}

In the first integral defining $E_P$, substitute $y=x-i$;
its contour is $\mathbb R-i$.  In the second substitute $y=x+i$;
its contour is $\mathbb R+i$.  The factors $e^{\pm t}$ disappear
under these substitutions.  Gamma recurrence gives exactly
\begin{equation}\label{gb:gamma-cancellation}
 (iy-\tfrac12)^{2m}W_m(y+i)
 =(iy+\tfrac12)^{2m}W_m(y-i)=2^{2m}V_m(y),
\end{equation}
where
\[
 V_m(y)=\left[
 \Gamma(\tfrac34+iy/2)\Gamma(\tfrac34-iy/2)\right]^{2m}.
\]
For example the first equality shifts
$\Gamma(-1/4+iy/2)$ to $\Gamma(3/4+iy/2)$; its recurrence
factor is $iy/2-1/4$.  The second equality has the opposite sign
before taking the even power, hence the same factor $2^{2m}$.

Every allowed multiplier satisfies $P(y+i)=P(y-i)$.  Both shifted
integrands are therefore the same meromorphic function
\[
 \mathcal F_P(y)=2^{2m}e^{iyt}V_m(y)P(y+i).
\]
The nearest Gamma poles of $V_m$ are $y=\pm3i/2$, outside
$|\Im y|\le1$.  For $P(x)=\sinh(\pi kx)$, $1\le k<m$,
the multiplier is $P(y+i)=(-1)^k\sinh(\pi ky)$, an entire
function.  Uniform Stirling estimates on the closed strip give
a polynomial times $e^{-\pi(m-k)|\Rea y|}$ as a bound for the
integrand, with a constant depending on the fixed real $t$.
For the tanh multiplier, $P(y+i)=\tanh(\pi y)$; away from its
finite poles it is bounded at the two ends of the strip, and
the decay is a polynomial times $e^{-m\pi|\Rea y|}$.
Thus the vertical sides of the rectangle between the two horizontal
contours vanish as their real parts tend to infinity.  Its bottom
side runs from left to right, so bottom minus top equals
$+2\pi i$ times the enclosed residue sum.

In the sinh case there are no enclosed poles and $E_P(t)=0$.
In the tanh case the only enclosed poles are the simple poles
$y=\pm i/2$.  Each tanh residue is $1/\pi$, and
$V_m(i/2)=V_m(-i/2)=\pi^m$.  Hence
\[
 \Res_{y=\pm i/2}\mathcal F_P(y)
      =2^{2m}\pi^{m-1}e^{\mp t/2},
\]
and
\begin{equation}\label{gb:residue-result}
 E_{\tanh}(t)=i\,2^{2m+2}\pi^m\cosh(t/2).
\end{equation}
All shifts have now been justified using pole-free Gamma factors,
the two explicitly listed tanh poles, and decaying vertical sides.

The second summand on the left of \eqref{gb:contour-algebra} is
analytic at zero.  Its first summand has Laurent coefficient
$2Q_P^{(2m)}(0)$ at $t^{-1}$.  Comparing this coefficient shows
that $Q_P^{(2m)}(0)=0$ for a sinh multiplier and
$Q_P^{(2m)}(0)=i2^{2m+2}\pi^m$ for tanh.  Equations
\eqref{gb:contour-algebra} and \eqref{gb:residue-result} therefore
give the exact germ identity
\[
 (Q_P^{(2m)}-Q_P^{(2m)}(0))\coth(t/2)
                  +(Q_P\coth(t/2))^{(2m)}=0.
\]
Finally set $R_P(z)=iQ_P(iz)$.  Pairing positive and negative
$x$ in $G_P(iz)$ gives precisely \eqref{gb:R}.  Rotation uses
$\coth(iz/2)=-i\cot(z/2)$ and yields
$\mathcal T_mR_P=0$.  Since
$R_P^{(2m)}(0)=i(-1)^mQ_P^{(2m)}(0)$, it also gives
\eqref{gb:top-derivative}.
\end{proof}

\begin{corollary}[Exact ordinary Gamma moments]\label{gb:moment-identities}
Put $\phi_n(x)=\Im(x+i/2)^{2n}$.  The Taylor coefficients of
\eqref{gb:R} are
\begin{equation}\label{gb:t}
 t_{n,j}:=[z^{2n}]R_j(z)
      =\frac4{(2n)!}\int_0^\infty\phi_n(x)P_j(x)W_m(x)\dd x.
\end{equation}
For every $m\ge1$,
\begin{align}
 \int_0^\infty\phi_m(x)\sinh(\pi kx)W_m(x)\dd x&=0
                  &&(1\le k<m),\label{gb:orthogonality}\\
 \int_0^\infty\phi_m(x)\tanh(\pi x)W_m(x)\dd x
          &=(-1)^{m+1}2^{2m}\pi^m.\label{gb:evaluation}
\end{align}
All these integrals converge absolutely.
\end{corollary}

\begin{proof}
The elementary generating identity
\[
 \sin(z/2)\sinh(xz)=
       \sum_{n\ge1}\frac{\phi_n(x)}{(2n)!}z^{2n}
\]
and local uniform exponential domination justify coefficient
extraction.  Now apply \eqref{gb:top-derivative}.
\end{proof}

The vanishing moments have a concrete polynomial interpretation.
Since $\phi_m(x)/x$ has degree $2m-2$ and its positive zeros are
\[
                  \frac12\cot\frac{r\pi}{2m},\qquad1\le r<m,
\]
equation \eqref{gb:orthogonality} is an explicit family of
multiple-orthogonality relations against positive weights on
$(0,\infty)$.  The zeros follow by writing
$x+i/2=|x+i/2|e^{i\theta}$, with $0<\theta<\pi/2$, and solving
$\sin(2m\theta)=0$.  These relations do not imply that every
later consecutive coefficient determinant has one fixed sign.

\subsection{A positive initial determinant and the exact Gamma representation}

Write
\[
 T_m=(t_{n,j})_{1\le n,j\le m},\qquad
 \Delta(y_1,\ldots,y_m)=\prod_{r<s}(y_s-y_r).
\]

\begin{proposition}[Positive normalization]\label{gb:normalization}
The initial determinant is
\begin{equation}\label{gb:detT}
                 \det T_m=(-1)^{m(m-1)/2}K_m,
\end{equation}
where $K_m>0$ has the absolutely convergent integral representation
\begin{equation}\label{gb:Km}
\begin{split}
 K_m={}&\frac{2^{(m-1)(m-2)/2}}{m!}
             \left(\prod_{n=1}^m\frac{4n}{(2n)!}\right)\\
 &\times\int_{(0,\infty)^m}
       \Delta(x_1^2,\ldots,x_m^2)
       \Delta(\cosh\pi x_1,\ldots,\cosh\pi x_m)\\
 &\hspace{38mm}\times
       \prod_{r=1}^m x_rW_m(x_r)\tanh(\pi x_r)
                    \dd x_1\cdots\dd x_m.
\end{split}
\end{equation}
Consequently the $R_j$ form a basis of the full formal null space,
and for every $b>m$,
\begin{equation}\label{gb:exact-Gamma-determinant}
 \boxed{\displaystyle
 \det F_{b,m}=
  \frac{(-1)^{m(m-1)/2}}{K_m}
        \det(t_{b+i,j})_{1\le i,j\le m}.}
\end{equation}
\end{proposition}

\begin{proof}
The polynomial $\phi_n(x)$ has leading term $nx^{2n-1}$ and
only lower odd powers thereafter.  A triangular row transformation
therefore gives
\[
 \det T_m=\left(\prod_{n=1}^m\frac{4n}{(2n)!}\right)
       \det\left(\int_0^\infty x^{2n-1}P_j(x)W_m(x)\dd x\right)_{n,j=1}^m.
\]
The continuous Cauchy--Binet identity writes the remaining determinant
as $1/m!$ times the integral of the two alternants
$\det(x_r^{2n-1})_{n,r}$ and $\det(P_j(x_r))_{r,j}$ against
$\prod_rW_m(x_r)\dd x_r$.  This identity follows directly by
expanding the two determinants, applying Fubini, and permuting the
integration variables; absolute convergence justifies each operation.
The first alternant is $\prod_rx_r\,\Delta(x_1^2,\ldots,x_m^2)$.

For the second alternant put $u=\pi x$, $v=\cosh u$, and use
\[
              \frac{\sinh(ku)}{\tanh u}=vU_{k-1}(v),
\]
where $U_{k-1}$ is the Chebyshev polynomial of the second kind,
with degree $k-1$ and leading coefficient $2^{k-1}$.
After factoring $\tanh(\pi x_r)$ from row $r$, the column
degrees are $m-1,m-2,\ldots,1,0$.  Their leading-coefficient
product is $2^{(m-1)(m-2)/2}$.  Reversing the columns thus gives
\[
\begin{split}
 \det(P_j(x_r))={}&(-1)^{m(m-1)/2}2^{(m-1)(m-2)/2}\\
 &\times\prod_r\tanh(\pi x_r)
              \Delta(\cosh\pi x_1,\ldots,\cosh\pi x_m).
\end{split}
\]
This proves \eqref{gb:detT}--\eqref{gb:Km}.  For distinct
positive $x_r$, the two Vandermondes have the same sign, since
both $x\mapsto x^2$ and $x\mapsto\cosh(\pi x)$ are strictly
increasing.  The integrand in \eqref{gb:Km} is strictly positive
away from the collision hyperplanes.  Its exponential growth from
the cosh Vandermonde is at most
$\prod_re^{(m-1)\pi x_r}$, while each $W_m(x_r)$ supplies
$e^{-m\pi x_r}$ times a polynomial factor.  At zero all factors
are bounded or vanish.  This proves convergence and $K_m>0$.

All $R_j$ are null germs, and their first $m$ coefficient columns
are independent by \eqref{gb:detT}.  The formal null space has
dimension $m$ by \eqref{gb:coefficient}, so they form a basis.
Their coefficient vectors obey \eqref{gb:recurrence}; consequently
\[
                  (t_{b+i,j})_{i,j=1}^m=F_{b,m}T_m.
\]
Taking determinants proves \eqref{gb:exact-Gamma-determinant}.
\end{proof}

\subsection{A Laplace-moment lemma with controlled error}

The asymptotic argument must retain one term beyond the leading
Laplace phase.  At every even mode that leading phase is real and
its imaginary part vanishes.  The following lemma supplies the
necessary next term without requiring an analytic continuation of
the error in a real density.

\begin{lemma}\label{gb:Laplace}
Let $m\ge1$ be fixed, $R>0$, and $a\in\mathbb R\setminus\{0\}$.
Suppose a real, locally bounded function $f$ on $[0,\infty)$ satisfies
\[
             f(x)=ax^{-m}e^{-Rx}(1+O(x^{-2}))\qquad(x\to\infty).
\]
As integers $N\to\infty$,
\begin{equation}\label{gb:Laplace-formula}
 \int_0^\infty(x+i/2)^Nf(x)\dd x
 =ae^{iR/2}\frac{\Gamma(N-m+1)}{R^{N-m+1}}
       \left(1+\frac{imR}{2(N-m)}+O(N^{-2})\right).
\end{equation}
The implicit constant may depend on $f,m,R$.
\end{lemma}

\begin{proof}
Choose fixed $X>1$ beyond which the asserted density estimate holds.
The integral over $[0,X]$ is $O(C^N)$ for a fixed $C$, hence
negligible compared with
$S_N=\Gamma(N-m+1)R^{-N+m-1}$ by Stirling's formula.
The absolute error made on $[X,\infty)$ by replacing $f(x)$ with
$ax^{-m}e^{-Rx}$ is bounded by a constant times
\[
             \int_X^\infty|x+i/2|^Nx^{-m-2}e^{-Rx}\dd x.
\]
Split this integral at $x=N^{2/3}$.  On the lower part it is at most
$\exp(\tfrac23N\log N+O(N))$, whereas
$S_N=\exp(N\log N+O(N))$.  This part is smaller than any fixed
inverse power of $N$ times $S_N$.  On the upper part,
\[
 \frac{|x+i/2|^N}{x^N}
   =\left(1+\frac1{4x^2}\right)^{N/2}
   \le\exp\!\left(\frac{N}{8x^2}\right)
\]
is uniformly bounded.  Its integral is therefore at most a constant
times $\Gamma(N-m-1)R^{-N+m+1}=O(N^{-2}S_N)$.
This estimate uses only real-variable bounds on the density error.

For the model integral set $y=x+i/2$.  After factoring $e^{iR/2}$,
its integrand is
$y^{N-m}(1-i/(2y))^{-m}e^{-Ry}$ on the ray
$[X,\infty)+i/2$.  Shift that ray to $[X,\infty)$ inside
$\Rea y\ge X$.  No pole is crossed.  The short vertical segment
at $X$ contributes $O(C^N)$, and the vertical segment at infinity
vanishes for each $N$ because of $e^{-R\Rea y}$.
Uniformly for real $y\ge X$,
\[
                 (1-i/(2y))^{-m}=1+\frac{im}{2y}+O(y^{-2}).
\]
Integration of the first two terms, extending their lower limit to
zero at another $O(C^N)$ cost, gives
\[
 ae^{iR/2}\left\{
 \frac{\Gamma(N-m+1)}{R^{N-m+1}}
 +\frac{im}{2}\frac{\Gamma(N-m)}{R^{N-m}}\right\}.
\]
The last term has absolute size $O(N^{-2}S_N)$ by the same Gamma
ratio estimate.  This proves \eqref{gb:Laplace-formula}.
\end{proof}

Define the constants
\begin{equation}\label{gb:mode-constants}
 \lambda_j=j\pi,\qquad \xi_j=(j\pi)^{-2},\qquad
 A_m^*=2^m(2\pi)^{2m},\qquad
 a_j=\begin{cases}A_m^*/2,&j<m,\\A_m^*,&j=m.
 \end{cases}
\end{equation}
Conjugate Gamma asymptotics contain only even inverse powers after
the leading factor.  Thus
\begin{equation}\label{gb:density-asymptotic}
 P_j(x)W_m(x)=a_jx^{-m}e^{-\lambda_jx}(1+O(x^{-2})).
\end{equation}
The further exponential terms in sinh or tanh are absorbed into the
indicated error.  In particular, there is no $x^{-1}$ correction.

\begin{proposition}[Distinct coefficient scales and exact radii]
\label{gb:coefficient-asymptotic}
Put $d_j=m+\ind_{2\mid j}$ and
\begin{equation}\label{gb:Lj}
 L_j=\begin{cases}
 4a_j\sin(j\pi/2)\lambda_j^{m-1},&j\text{ odd},\\
 2ma_j\cos(j\pi/2)\lambda_j^m,&j\text{ even}.
 \end{cases}
\end{equation}
All $L_j$ are nonzero, and
\begin{equation}\label{gb:t-asymptotic}
 t_{n,j}=L_j\lambda_j^{-2n}(2n)^{-d_j}(1+O(n^{-1})).
\end{equation}
The Taylor radius of $R_j$ is exactly $\lambda_j=j\pi$.
\end{proposition}

\begin{proof}
Apply \cref{gb:Laplace} to \eqref{gb:density-asymptotic}, take
the imaginary part, and multiply by $4/(2n)!$.  If $j$ is odd,
$e^{i\lambda_j/2}=\pm i$ and the leading term survives.  Since
$\Gamma(2n-m+1)/\Gamma(2n+1)\sim(2n)^{-m}$, this gives the
first line of \eqref{gb:Lj}.  If $j$ is even, the leading phase
is $\pm1$ and its imaginary part is zero.  The next term in
\eqref{gb:Laplace-formula} gives instead
\[
 \frac{4a_j\cos(j\pi/2)}{(2n)!}
 \frac{m}{2}\frac{\Gamma(2n-m)}{\lambda_j^{2n-m}},
\]
which is the second line of \eqref{gb:Lj} and the exponent $m+1$.
The $O(N^{-2})$ error in the lemma is small enough to give a
relative $O(n^{-1})$ error even in this second case.  Finally
\eqref{gb:t-asymptotic} implies
$\lim_{n\to\infty}|t_{n,j}|^{1/(2n)}=\lambda_j^{-1}$.
The Cauchy--Hadamard formula proves the exact radius.  This also
provides another proof of independence: the smallest surviving
radius in any nontrivial linear combination cannot cancel.
\end{proof}

\subsection{Eventual optimal reconstruction for every fixed gap}

Set
\begin{equation}\label{gb:asymptotic-normalization}
 D_m=m^2+\lfloor m/2\rfloor,\qquad
 \mathcal A_m=
 \frac{\prod_{j=1}^m|L_j|}{K_m}
       \left(\prod_{j=1}^m\xi_j\right)
       \prod_{j<k}(\xi_j-\xi_k)>0.
\end{equation}

\begin{theorem}[Fixed-gap determinant asymptotic]\label{gb:eventual}
For each fixed integer $m\ge1$, as $b\to\infty$,
\begin{equation}\label{gb:det-asymptotic}
 \boxed{\displaystyle
 \det F_{b,m}\sim
 (-1)^{\lfloor m^2/4\rfloor}\mathcal A_m
 (2b)^{-D_m}\bigl[\pi^{2m}(m!)^2\bigr]^{-b}.}
\end{equation}
In particular, there exists an integer $B(m)>m$ such that every
pair $(p,q)=(2b-2m,2b)$ with $b\ge B(m)$ has an optimal first-$q$
formal moment prefix.  Its late-block determinant has eventual sign
\begin{equation}\label{gb:eventual-sign}
 \operatorname{sgn}\det M_{b-m,b}
        =(-1)^{b(b-1)/2+\lfloor m^2/4\rfloor}.
\end{equation}
\end{theorem}

\begin{proof}
From column $j$ of $(t_{b+i,j})_{i,j=1}^m$, factor
$L_j\lambda_j^{-2b}(2b)^{-d_j}$.  By
\eqref{gb:t-asymptotic}, the remaining matrix tends entrywise to
$V=(\xi_j^i)_{i,j=1}^m$.  Its determinant is
\[
 \det V=\left(\prod_j\xi_j\right)
                   \prod_{j<k}(\xi_k-\xi_j)\ne0.
\]
The distinct numbers $\xi_1>\cdots>\xi_m>0$ give sign
$(-1)^{m(m-1)/2}$, which cancels the same sign in
\eqref{gb:exact-Gamma-determinant}.  Furthermore
\[
 \prod_j\operatorname{sgn}L_j
    =(-1)^{\sum_{j=1}^m\lfloor j/2\rfloor}
    =(-1)^{\lfloor m^2/4\rfloor},\qquad
 \sum_jd_j=m^2+\lfloor m/2\rfloor,
\]
and
$\prod_j\lambda_j^{-2b}=[\pi^{2m}(m!)^2]^{-b}$.
These identities prove \eqref{gb:det-asymptotic} with the positive
constant \eqref{gb:asymptotic-normalization}.  A sequence with a
nonzero asymptotic equivalent is nonzero and has the same sign
for all sufficiently large indices.  Apply the exact positive-factor
bridge \eqref{gb:bridge-formula} and then \eqref{gb:full-prefix}.
\end{proof}

\begin{remark}[Scope of the theorem]\label{gb:scope}
The assertion is eventual for each \emph{fixed gap} $2m$ and
does not provide a numerical value of $B(m)$.  It leaves only
finitely many possible exceptional pairs on each such ray, but
does not prove the entire gap-six family or the general
all-unequal-even conjecture.  No uniform assertion as $m$ grows
with $b$ follows from the fixed-$m$ errors used here.
\end{remark}

For $m=1$, \eqref{gb:evaluation} gives $K_1=8\pi$, and
\eqref{gb:asymptotic-normalization} gives $\mathcal A_1=4/\pi$.
The fundamental adjacent-gap sequence therefore satisfies
\[
        C^{[1]}_{b+1,1}\sim\frac{2}{\pi b}\,\pi^{-2b}.
\]
For $m=2$ the eventual sign of $\det F_{b,2}$ is negative,
consistent with the inherited theorem for every gap-four pair.
For $m=3$ it is positive.  Finite exact checks give
$\det F_{b,3}<0$ at $b=4,5,6$ and $\det F_{b,3}>0$ at
$7\le b\le40$.  These observations illustrate a sign transition;
the proof of eventual positivity is \cref{gb:eventual}, and the
finite checks are not a proof of nonvanishing for every gap-six pair.

\subsection{All-orders coefficients and a harmonic-number correction}

The Gamma construction supplies more than a leading equivalent.
It gives an explicit Poincare expansion to every finite order,
and the first determinant correction has a concise harmonic-number
form.  Let Euler numbers be defined by
\[
                \operatorname{sech}z=
                \sum_{k\ge0}E_{2k}\frac{z^{2k}}{(2k)!}.
\]
Define rational numbers $w_{m,r}$ by the formal identity
\begin{equation}\label{gb:w-series}
 \sum_{r\ge0}w_{m,r}X^r
   =\exp\!\left(\sum_{k\ge1}
            \frac{m|E_{2k}|}{k\,2^{2k+1}}X^k\right).
\end{equation}
Thus $w_{m,0}=1$ and $w_{m,1}=m/8$.
For $\ell\ge0$ set
\begin{equation}\label{gb:b-coefficients}
 b_{m,\ell}=\sum_{r=0}^{\lfloor\ell/2\rfloor}
  w_{m,r}\frac{(m+2r)_{\ell-2r}}{(\ell-2r)!}
                     (i/2)^{\ell-2r}.
\end{equation}

\begin{proposition}[All-orders mode construction]\label{gb:all-orders}
For each fixed mode $j$, with $\lambda_j,a_j$ as in
\eqref{gb:mode-constants},
\begin{equation}\label{gb:all-orders-moments}
 \int_0^\infty(x+i/2)^NP_j(x)W_m(x)\dd x
 \sim a_je^{i\lambda_j/2}
       \sum_{\ell\ge0}b_{m,\ell}
       \frac{\Gamma(N-m-\ell+1)}{
                            \lambda_j^{N-m-\ell+1}}.
\end{equation}
This is a Poincare expansion: truncation after any fixed $\ell=L$
has error of size
$O(\Gamma(N-m-L)\lambda_j^{-N+m+L})$.
Taking imaginary parts retains even $\ell$ when $j$ is odd and
odd $\ell$ when $j$ is even.  In particular,
\begin{equation}\label{gb:t-first-correction}
 t_{n,j}=L_j\lambda_j^{-2n}(2n)^{-d_j}
       \left(1+\frac{d_j(d_j-1)}{4n}+O(n^{-2})\right).
\end{equation}
\end{proposition}

\begin{proof}
The shifted logarithmic Stirling expansion is
\[
 \log\Gamma(z+\alpha)\sim
 (z+\alpha-\tfrac12)\log z-z+\tfrac12\log(2\pi)
 +\sum_{n\ge1}\frac{(-1)^{n+1}B_{n+1}(\alpha)}{n(n+1)z^n},
\]
with a remainder of the next prescribed algebraic order in each
closed subsector of $|\arg z|<\pi$ \cite{dlmf-stirling}.
Apply it with $\alpha=1/4$ and $z=\pm ix/2$.
The conjugate pair cancels odd inverse powers.  Using
\[
 B_{2k+1}(1/4)=-\frac{(2k+1)E_{2k}}{4^{2k+1}}
\]
gives
\[
 \log W_m(x)\sim
 \log A_m^*-m\log x-m\pi x
       +\sum_{k\ge1}\frac{m|E_{2k}|}{k\,2^{2k+1}x^{2k}}.
\]
Exponentiation to any fixed order proves
\begin{equation}\label{gb:W-all-orders}
 W_m(x)\sim A_m^*x^{-m}e^{-m\pi x}
                       \sum_{r\ge0}w_{m,r}x^{-2r}.
\end{equation}
The corrections to the leading exponential in $P_j$ are themselves
exponentially smaller, so they contribute less than any fixed
inverse power after the leading density has been removed.

For each term $x^{-m-2r}e^{-\lambda_jx}$ in the resulting density,
repeat the shift $y=x+i/2$ from \cref{gb:Laplace}.  Its algebraic
factor becomes
\[
 y^{-m-2r}\left(1-\frac{i}{2y}\right)^{-m-2r}
   =y^{-m-2r}\sum_{h\ge0}
       \frac{(m+2r)_h}{h!}\left(\frac{i}{2y}\right)^h.
\]
Collecting terms with $2r+h=\ell$ gives exactly
\eqref{gb:b-coefficients}.  To justify any fixed truncation, split
at fixed $X>1$ as in the lemma, retain the finitely many necessary
density terms, and bound its real remainder by the first omitted
power of $x^{-1}$ times the leading exponential.  The split at
$N^{2/3}$ bounds its moment by the corresponding Gamma ratio.
The finitely many model contours have the same pole-free deformation,
and Taylor's remainder for $(1-i/(2y))^{-m-2r}$ is uniform for
$y\ge X$.  These bounds prove the stated remainder in
\eqref{gb:all-orders-moments}; no contour continuation of the real
Stirling remainder is used.

For even $\ell$ the coefficient $b_{m,\ell}$ is real, and for
odd $\ell$ it is purely imaginary.  Thus the phase
$e^{ij\pi/2}$ gives the asserted parity selection.  At odd $j$
the leading imaginary term has $\ell=0$ and the next possible
one has $\ell=2$.  At even $j$ the leading one has $\ell=1$,
with $b_{m,1}=im/2$, and the next possible one has $\ell=3$.
In either case the first relative inverse-power correction comes
solely from the factorial ratio
\[
 \frac{\Gamma(N-d+1)}{\Gamma(N+1)}
 =N^{-d}\left(1+\frac{d(d-1)}{2N}+O(N^{-2})\right)
 \qquad(d\in\mathbb N).
\]
For integer $d$ this follows directly by expanding
$\prod_{r=0}^{d-1}(N-r)^{-1}$.  Set $N=2n$ and $d=d_j$
to obtain \eqref{gb:t-first-correction}.
\end{proof}

\begin{theorem}[The harmonic first correction]\label{gb:harmonic-correction}
Let $e=\lfloor m/2\rfloor$ and $H_0=0$.  Define
\begin{equation}\label{gb:kappa}
 \boxed{\displaystyle
 \kappa_m=-\frac{m^2(m+3)+e(2m+3)}4
       +\sum_{r=1}^{e}r\bigl(H_{m+2r}-H_{m-2r}\bigr).}
\end{equation}
Then
\begin{equation}\label{gb:det-first-correction}
\begin{split}
 \det F_{b,m}={}&(-1)^{\lfloor m^2/4\rfloor}\mathcal A_m
       (2b)^{-D_m}\bigl[\pi^{2m}(m!)^2\bigr]^{-b}\\
 &\times\left(1+\frac{\kappa_m}{b}+O(b^{-2})\right).
\end{split}
\end{equation}
More generally, this determinant has an expansion to every finite
inverse-power order.  For fixed $m$, all its normalized coefficients
belong to $\Q[\pi^2]$.
\end{theorem}

\begin{proof}
Factor the same column scales as in the proof of
\cref{gb:eventual}.  By \eqref{gb:t-first-correction}, the remaining
entry in row $i$, column $j$ is
\[
 \xi_j^i\left\{1+\frac1b
           \left(\frac{d_j(d_j-1)}4-id_j\right)+O(b^{-2})\right\}.
\]
For $V=(\xi_j^i)_{i,j=1}^m$, the first relative correction in its
determinant is therefore
\begin{equation}\label{gb:kappa-trace}
 \frac14\sum_{j=1}^md_j(d_j-1)-\sum_{j=1}^md_j\eta_j,
 \qquad
 \eta_j=\left(V^{-1}\operatorname{diag}(1,\ldots,m)V\right)_{jj}.
\end{equation}
The diagonal quantities can be computed without matrix inversion.
Differentiating the Vandermonde product gives
\[
 \eta_j=\xi_j\frac{\partial}{\partial\xi_j}\log\det V
       =1+\sum_{k\ne j}\frac{\xi_j}{\xi_j-\xi_k}.
\]
For $\xi_j=(j\pi)^{-2}$, partial fractions give
\begin{equation}\label{gb:eta-harmonic}
 \eta_j=m+\frac34+\frac j2(H_{m-j}-H_{m+j}).
\end{equation}
Indeed $\xi_j/(\xi_j-\xi_k)=k^2/(k^2-j^2)$, and
\[
 \sum_{k\ne j}\left(\frac1{k-j}-\frac1{k+j}\right)
        =H_{m-j}-H_{m+j}+\frac{3}{2j}.
\]
Also $\sum_j\eta_j=m(m+1)/2$, either by the trace or by
pairing the summands in the Vandermonde expression.  Since
$d_j=m+\ind_{2\mid j}$, substitution into
\eqref{gb:kappa-trace} gives
\[
 -\frac{m^2(m+3)}4+\frac{me}{2}
                         -\sum_{r=1}^e\eta_{2r}.
\]
Equation \eqref{gb:eta-harmonic} now simplifies it to
\eqref{gb:kappa}, proving \eqref{gb:det-first-correction}.

For higher orders, insert any desired finite truncation of
\eqref{gb:all-orders-moments} into the fixed-size determinant.
Its leading Vandermonde is nonzero, so division by that leading
term gives a valid Poincare expansion.  Relative to $L_j$, parity
selection leaves only even powers of $\lambda_j=j\pi$ in the
mode corrections.  The factorial-ratio expansions have rational
coefficients, while ratios of Vandermonde expressions in the
$\xi_j$ cancel their common powers of $\pi^{-2}$.  Thus each
normalized determinant coefficient is a polynomial in $\pi^2$
with rational coefficients, as claimed.
\end{proof}

For orientation, the first correction values are
\[
 \kappa_1=-1,\qquad\kappa_2=-\frac{14}{3},\qquad
 \kappa_3=-\frac{217}{15},\qquad\kappa_4=-\frac{949}{35}.
\]
The series in \eqref{gb:w-series} and
\eqref{gb:all-orders-moments} are formal asymptotic constructions;
convergence is not asserted.  The argument also does not establish
a complete exponentially improved expansion or an explicit error
constant sufficient to certify a numerical $B(m)$.

\subsection{A hypergeometric form of the null equation}

The differential equation itself has a useful classical
hypergeometric form.  Let $Q$ be any even analytic null germ of
$\mathcal L_m$ and write $c=Q^{(2m)}(0)$.  On the local branches
near $u=1$ with $\log1=0$ and $\sqrt1=1$, define
\[
 Y(u)=\frac{Q(\tfrac12\log u)}{\sqrt u-1},\qquad
                    \Theta=u\frac{d}{du}.
\]
The apparent singularity at $u=1$ is removable because $Q(t)=O(t^2)$.
Then
\begin{equation}\label{gb:hypergeometric}
 \left[\Theta^{2m}-u(\Theta+\tfrac12)^{2m}\right]Y(u)
            =-\frac{c}{2^{2m+1}}(1+\sqrt u).
\end{equation}
To verify it, set $Q(t)=(e^t-1)F(t)$, so that
$Q(t)\coth(t/2)=(e^t+1)F(t)$.  Multiplication of
$\mathcal L_mQ=0$ by $e^t-1$ gives
\[
 D^{2m}F-e^{2t}(D+1)^{2m}F=-\tfrac c2(1+e^t).
\]
Now $u=e^{2t}$ and $D=2\Theta$, proving
\eqref{gb:hypergeometric}.  Its homogeneous operator is the standard
generalized hypergeometric operator associated with
\[
 {}_{2m}F_{2m-1}
 \left(\begin{matrix}\tfrac12,\ldots,\tfrac12\\1,\ldots,1\end{matrix};u\right)
\]
\cite{dlmf-hypergeometric}.  This operator identification is classical
algebra.  The Gamma basis selects and normalizes the relevant
analytic null germs at $u=1$ and connects them to the harmonic-tail
coefficient matrix; that connection is the substantive result here.

\subsection{Verification and the next finite obstruction}

The accompanying scripts separate exact algebra from numerical
diagnostics.  Direct Gamma quadrature at $45$ decimal digits for
$m=1,2,3$ and $1\le n\le m+3$ gives all 18 coefficient recurrence
residuals below $10^{-46}$ and all six top-derivative identity
residuals below $10^{-42}$, with the normalizations recorded
separately. A separate calculation evaluates $T_m$
by Gamma quadrature at $60$ digits, then propagates the exact
recurrence to compare the determinant with
\eqref{gb:det-asymptotic}.  The recorded ratios are
\begin{center}
\begin{tabular}{@{}crrrr@{}}
\toprule
 & $b=20$ & $b=40$ & $b=80$ & $b=120$\\
\midrule
$m=1$ & $0.951663$ & $0.975424$ & $0.987607$ & $0.991714$\\
$m=2$ & $0.796107$ & $0.891056$ & $0.943647$ & $0.961999$\\
$m=3$ & $0.456866$ & $0.687818$ & $0.832149$ & $0.885309$\\
\bottomrule
\end{tabular}
\end{center}
Their approach to one is a consistency check on the explicit
normalization; the infinite theorem is proved above independently
of these finite computations.

The most direct continuation is to make the remainder in
\eqref{gb:det-asymptotic} effective, obtain a certified cutoff for
each fixed $m$, and use exact rational arithmetic below it.  The
gap-six case is a natural first target because its initial sign
transition rules out a blanket positivity shortcut.  Further questions
include exact evaluations of $K_m$, a classification of finite
sign changes, and uniform asymptotics when $m$ grows with $b$.
The Gamma moments also suggest parameter deformations of
\eqref{gb:W} and \eqref{gb:cot-operator}; any such extension must
recheck the contour residues and the original coordinate bridge.
